% !PRINTED 363
\weberpage{363}

%<<CONT
one represents the ideals, according to Bd.~II, \S\,169, by the \gesperrt{functionals} discussed in Bd.~II, \S\,153\,ff., because in doing so one can make use of the ordinary rules of multiplication and division\footnote{Cf. also the author's memoir \glqq Über Zahlengruppen in algebraischen Körpern\grqq{}, first communication, Mathematische Annalen 48, 438.}.

\gesperrt{We restrict ourselves here once and for all, in the field $\Omega$, to numbers and ideals that are coprime to an arbitrarily assumed rational integer $Q$} (the conductor of an order)\gesperrt{.}

Under this supposition the integral and fractional numbers of the field $\Omega$ form the group previously denoted by $O$. We extend this group by the addition of the \gesperrt{functionals} and denote the \gesperrt{group} so extended by $\overline{O}$. To every element of $\overline{O}$ there then corresponds an integral or fractional ideal, and to those elements whose quotient is a functional unit \gesperrt{there corresponds the same ideal.}

If $\eta$ is a number in $O$ and $\varepsilon$ is any functional unit, then a \gesperrt{principal ideal} corresponds to the product $\varepsilon\eta$. Accordingly we denote
\begin{quote}
\gesperrt{by $\overline{E}$ the group of functional units,}\\
\gesperrt{\glqq\ $E$\ \glqq\ \glqq\ \glqq\ numerical units,}
\end{quote}
and by
\begin{quote}
\gesperrt{$\overline{E}O$ the principal class of the functionals.}
\end{quote}
\gesperrt{If we take a system of representatives}
\weq{1}{\varphi_1,\ \varphi_2,\ \dots,\ \varphi_h}
of non-equivalent integral functionals in $\Omega$, then every integral or fractional functional $\Phi$ will be representable in the form:
\[
\Phi = \varphi_i \overline{\omega},
\]
where $\overline{\omega}$ is a functional from $\overline{E}O$, and the classes of ideals in $\Omega$ are the subsidiary groups of $\overline{E}O$ in $\overline{O}$. Accordingly the class number $h$ is the index of the divisor $\overline{E}O$ of $\overline{O}$:
\weq{2}{h = (\overline{O},\ EO).} % [sic] German source prints E here without an overbar.
Here, in the sharper division into classes, one has to restrict oneself in the case of $O$ to the numbers with positive norm.

If we now consider the division of the ideals into classes according to the order $[Q]$, then we have to take the functional group $\overline{E}O'$ as the \gesperrt{principal class}, and the class number is
\[
h' = (\overline{O},\ \overline{E}O').
\]
